\documentclass{article}
\usepackage{style}

\title{Experiment 7 - The Hall Effect}
\author{Daire O'Brien - 24365676}
\begin{document}
\maketitle
    \section{Background}
        \subsection{Discovery}
            In 1879, based on work by André-Marie Ampére and James Maxwell, Edwin Hall investigated how magnetic fields interacted with an electrical current in a conductor. He proposed that 
                \[E' \propto Mv\]
            where $E' \equiv V$, $M \equiv \abs{\vec B}$ and "$v$ is the velocity of the electricity in the [object]" (E. H. Hall, 1879)$^{\text{[1]}}$ (At the time of Hall's paper the electron had only been theorized and wasn't discovered until ) 
                \[V_H = \frac {I_x B_z}{n_e C_e t}\]
            where $I_x$ is the current through the conductor, $B_z$ is the magnetic field, $n_e$ is the number of electrons, $C_e$ is the charge of the electron and $t$ is the thickness of the material.

        \subsection{Hall Effect Sensors}
            Hall effect sensors typically combine a Hall element and an op-amp. The Hall element has a similar structure to a Wheatstone bridge
            Hall effect sensors are often used to measure the position or angle of something, for example a stick in a game controller, the rotation speed or postion of a motor's shaft or measuring current without needing to be inserted directly into a given circuit.
    \section{References}
        \begin{enumerate}
            \item Edwin Hall, 1879, \emph{On a New Action of the Magnet on Electric Currents}, American Journal of Mathematics Vol. 2 No. 3
            \item Manny Soltero, 2023, \emph{What Is a Hall-effect Sensor?}, Texas Instruments technical article
        \end{enumerate}
\end{document}